\section{Discussion and Limitations}
\label{sec:discussion}

The central question posed by this work is not whether multi-agent reinforcement learning
can control traffic signals---this is well established---but how much of that capability
survives when the policy is denied every privileged simulator quantity and restricted to
signals a roadside camera can actually compute. The experiments answer this quantitatively.
The information-restriction cost, defined as the performance gap between the proxy policy and
the privileged upper bound under identical pristine features (Section~\ref{sec:restriction_cost}),
is only $\RestrictionCost\%$ in mean travel time, while the proxy policy itself improves on the strongest
classical baseline by $\ImproveStrongest\%$. The interpretation is direct: for the signal-timing objective,
the camera-observable feature set retains the overwhelming majority of information, and the
residual loss from discarding exact per-vehicle delays and network-wide counts is small
relative to the gains from learning itself. This isolates a result that prior camera-state RL
leaves implicit, because those methods constrain the observation while continuing to train on
a privileged reward~\cite{azfar2025_cosim,dai2022_worldmodels}; here the reward is itself
camera-computable, so the reported cost reflects a genuinely deployable controller rather than
a partially privileged one.

The mechanism behind this performance is illuminated by the ablations
(Section~\ref{sec:ablations}). The two reward components designed against starvation---the
mean-of-squared-queue penalty and the signed phase-imbalance term---contribute disproportionately to
the tail-latency metric: removing either inflates the P95 waiting time far more than the mean,
confirming that the quadratic penalty and the directional phase-imbalance signal act precisely where
a linear queue average and a one-sided pressure surrogate~\cite{varaiya2013_maxpressure,
wei2019_presslight} are blind. The advantage of centralized training over the independent-critic
variant~\cite{dewitt2020_ippo} further indicates that the observed gains are not merely the
product of a richer state, but of coordinated credit assignment across multiple intersections
under shared, non-stationary demand.

These findings should be read as closing an \emph{observation} gap, which is distinct from,
and complementary to, the \emph{dynamics} gap that dominates sim-to-real transfer in robotics
and in prior transfer-oriented TSC~\cite{peng2018_dynrand,zhao2020_sim2real_survey,
da2024_promptgat}. The graceful degradation observed across the $\{0\times,0.5\times,1\times,2\times\}$
sensing-noise envelope (Section~\ref{sec:robustness_noise}) shows that the policy does not rely
on brittle, high-precision feature recovery: performance declines smoothly rather than
collapsing as perception error grows, and the relative ordering against the heuristic baselines
is preserved throughout. We emphasize that this is evidence about the policy's tolerance to
bounded perception error under a documented-basis noise model, not a claim of guaranteed field
transfer; the latter requires the dynamics gap to be addressed in tandem and is the subject of
the companion study.

From a practitioner's standpoint, the practical implication is that adaptive, learned signal
control need not wait for instrumented intersections or privileged feedback channels. A single
camera and a legacy two-phase controller---the prevailing hardware in many developing
metropolises---are sufficient to host the policy, since every quantity required at inference
is derived from detection and tracking alone. This lowers the deployment barrier from a
sensing-infrastructure problem to a calibration-and-integration problem.

\noindent\textbf{Limitations and future work.}
While this study isolates the information-restriction cost of camera-computable control,
several limitations remain. First, the evaluation employs a two-phase logic without protected
turns, matching legacy hardware but limiting international generality. Second, the 1D effective
queue is a scalar proxy for 2D sublane flow, and the local throughput reward introduces a minor
theoretical risk of downstream spillback (greedy flushing) that the central critic only
partially mitigates. Finally, the findings establish the robustness boundary in simulation;
bridging the physical dynamics gap with measured feature-recovery calibration is left to the
companion field deployment.
